\section{AV1: a local state bound linear in the coupling}
\label{sec:av1}

\textbf{Contribution \projecttag{HNM-C-AV1}.} The AV1 contract, recorded under
the title ``Hruday first-order local state lemma for the zero-selected AQ
subfamily'', was frozen before production at the unchanged cap
$\tau=\pm10^{-8}$. It preregistered the target $D_{\rm ii}\le4/10^7$ for the
exact-first-order tier. That is the value AV2 needs, since
$2(D+D^2)+49\cdot10^{-8}/\pi\le10^{-6}$ requires $D\le4.22\times10^{-7}$; the
contract kept $10^{-6}$ as a secondary comparison. Two routes were assigned.
The forward producer derives the reduced density explicitly. The reverse
producer, under premise isolation, bounds only the vacuum overlap. The
product-ordering split and the location of the normalization trap are shared
panel premises. Independence is claimed only for the inequality, the
constants, the enumeration, the cutoff step and the passage step.

\subsection{Reviewed scope}
The \repo{research/round32/advisor/av1-gate.json}{AV1 gate} records the
verdict \recid{accepted_within_scope}. Its accepted statement, quoted verbatim
below in the gate's ASCII notation, is the scope of every claim in this
section.
\begin{gatequote}
In the AM2/AQ1 zero-selected patterned family (SU(2) Kogut-Susskind form on
Z\textasciicircum{}3 at fixed spacing, coarse 24-link factors, selected triple
exactly (0,0,0) with Haar product reference, 21 omitted anchored faces per
factor grouped into whole stars phi\_b=-(tau/3) sum W\_f in normalized units
delta=alpha/8), at either sign of tau with |tau|<=10\textasciicircum{}-8 (the
cap values bound every smaller |tau| because each bound is increasing in
|tau|), the unique untruncated ground state of every centered whole-star box
Lambda\_N, N>=2, has reduced density on the complete cover R=\{0,e\_z\} of the
original xz Wilson loop satisfying ||rho\_\{N,R\}-P\_R||\_1<=D uniformly in N.
The bound follows from the product-ordering split psi=psi\_out+delta, (P\_R x
1)psi=psi\_out, ||delta||<=eps||psi\_out|| (straddling supports and the
two-creation term included) by two independently proved inequalities: the
explicit-density bound 2eps(1+eps)/(1+eps\textasciicircum{}2) (forward) and
the fidelity bound 2eps/sqrt(1+eps\textasciicircum{}2) (reverse). Tier (ii),
with the exact first-order face coefficient |tau|/144 and the self-consistent
AM2 remainder t<=t\_1/(1-352J):
D\_ii=585079838465912592144137406066050/42981220507576537932303142777593983768257
(\textasciitilde{}1.3612e-8; exact rational, certified by both inequalities),
and the reverse\textquotesingle{}s labelled 82-face refinement
D\_ii<=113902305553947264976096174312887/10\textasciicircum{}40
(\textasciitilde{}1.1390e-8; directed rational upper bound via the fidelity
inequality); both meet 4/10\textasciicircum{}7 and 10\textasciicircum{}-6 at
both signs. Tier (i), crude majorant t<=J\_0G(R)=37/6250000:
D\_i=36133347268157653748322/1525878906463907516326874161
(\textasciitilde{}2.3680e-5; exact, forward) and
D\_i<=236800700911401877571290493559933441/10\textasciicircum{}40
(\textasciitilde{}2.3680e-5; directed, reverse), failing both targets and
retained. Both tiers scale linearly (D(tau)/D(tau/100) in [99,101]).
Consequences: |omega(W)|<=D, |omega(W\textasciicircum{}2)-1/4|<=D/2, with
m\textasciicircum{}2<=D\textasciicircum{}2 charged. The on-site cutoff is
removed for the ground vector itself in each fixed box (uniform-gap argument),
and the bound passes by local trace-norm convergence to every subsequential
limit of AQ1\textquotesingle{}s centered whole-star construction, including AQ1\textquotesingle{}s chosen state:
uniform local closeness, not uniqueness, whole-sequence convergence or a rate
in N. No Euclidean node certificate, interaction-induced shift or its sign,
first-order coefficient, uniform Wilson, weak-coupling or continuum claim
follows.
\end{gatequote}
The rest of this section derives that statement. The limitations recorded by
the gate close the section (Section~\ref{sec:av1-limits}).

\subsection{Product ordering and the exact split}
\label{sec:av1-split}
Fix a centered whole-star box $\Lambda_N$, $N\ge2$, and let $\psi=e^{-C}\Omega_0$
be the AM2 ground vector of \eqref{eq:dict-fixed}. Because the creations
commute and square to zero, $e^{-C}$ factorizes, and the factors whose supports
meet $R$ can be placed last:
\begin{equation}
 e^{-C}=\prod_I(1-\hat c_I),\qquad
 \psi=\prod_{I\cap R\ne\varnothing}(1-\hat c_I)\,\psi_{\rm out},\qquad
 \psi_{\rm out}=\prod_{I\cap R=\varnothing}(1-\hat c_I)\,\Omega_0
 =\Omega_R\otimes\varphi_{\rm out}.
 \label{eq:av1-order}
\end{equation}
The factorization of $\psi_{\rm out}$ holds because every creation disjoint
from $R$ acts as the identity on $\mathcal H_R$. In the remaining product, any
two creations with overlapping supports multiply to zero. Pairwise disjoint
supports that each meet the two-site set $R$ number at most two: one contains
$0$ but not $e_z$, the other $e_z$ but not $0$. Hence
\begin{equation}
 \delta:=\psi-\psi_{\rm out}
 =-\sum_{I\cap R\ne\varnothing}\hat c_I\,\psi_{\rm out}
 +\sum_{\substack{I\ni0,\ J\ni e_z\\ I\cap J=\varnothing}}\hat c_I\hat c_J\,\psi_{\rm out}.
 \label{eq:av1-delta}
\end{equation}
The second sum is the two-creation term. In the reverse form, with $C_R$ the
sum of creations meeting $R$, $C_R^3=0$ and
$\delta=-C_R\psi_{\rm out}+\tfrac12C_R^2\psi_{\rm out}$.

Every term of $\delta$ contains a creation that is excited at some $x\in R$,
and $P_xQ_x=0$. Therefore
\begin{equation}
 (P_R\otimes1)\psi=\psi_{\rm out},\qquad
 \langle\psi_{\rm out},\delta\rangle=0,\qquad
 \norm{\psi}^2=\norm{\psi_{\rm out}}^2+\norm{\delta}^2 .
 \label{eq:av1-orth}
\end{equation}
For a creation meeting $R$,
$\hat c_I\psi_{\rm out}=c_I\otimes\Omega_{R\setminus I}\otimes
(\langle\Omega_{I\setminus R}|\otimes1)\varphi_{\rm out}$, and the partial
vacuum bra is a contraction. The sum includes \emph{straddling} supports,
which meet both $R$ and its complement. With $t=\norm{c}_a$ and
$\norm{\psi_{\rm out}}=\norm{\varphi_{\rm out}}$,
\begin{equation}
 e:=\frac{\norm{\delta}}{\norm{\psi_{\rm out}}}\le\varepsilon
 :=\sum_{I\cap R\ne\varnothing}\norm{c_I}
 +\sum_{\substack{I\ni0,\ J\ni e_z\\ I\cap J=\varnothing}}\norm{c_I}\norm{c_J}
 \le2t+t^2 .
 \label{eq:av1-eps}
\end{equation}
The anchored maximum enters only at the two sites of $R$. A per-site maximum
summed over the volume would be extensive; it is a rejected control.

\subsection{Two inequalities}
\label{sec:av1-inequalities}
\paragraph{Forward: explicit reduced density.}
Put $n=\norm{\psi_{\rm out}}$, $d=\norm{\delta}$,
$\xi=(1_R\otimes\langle\varphi_{\rm out}|)\delta$ and
$\sigma=\tr_{\rm out}|\delta\rangle\langle\delta|$. Direct partial traces and
\eqref{eq:av1-orth} give, exactly,
\begin{equation}
 \rho_{N,R}=\frac{n^2P_R+|\Omega_R\rangle\langle\xi|+|\xi\rangle\langle\Omega_R|+\sigma}{n^2+d^2},
 \qquad \langle\Omega_R,\xi\rangle=0,\quad P_R\sigma=0,\quad \tr\sigma=d^2 .
 \label{eq:av1-density}
\end{equation}
Cauchy--Schwarz gives $\norm{\xi}\le nd$, and on $\mathrm{span}\{\Omega_R,\xi\}$
the cross operator has trace norm $2\norm{\xi}$. Hence
\begin{equation}
 \norm{\rho_{N,R}-P_R}_1\le\frac{2nd+2d^2}{n^2+d^2}=\frac{2e(1+e)}{1+e^2}
 \le f(\varepsilon):=\frac{2\varepsilon(1+\varepsilon)}{1+\varepsilon^2}.
 \label{eq:av1-forward}
\end{equation}
The function $f$ increases on $[0,1+\sqrt2]$, and $f(\varepsilon)\ge2$ once
$\varepsilon\ge1$, so the last step holds for every $\varepsilon$.

\paragraph{Reverse: vacuum overlap and fidelity.}
By \eqref{eq:av1-orth}, $\tr(\rho_{N,R}P_R)=n^2/(n^2+d^2)=1/(1+e^2)$ exactly.
The pure-state mixture inequality admitted in AT4,
$\norm{\rho-|\Omega\rangle\langle\Omega|}_1\le2\sqrt{1-\langle\Omega,\rho\,\Omega\rangle}$,
then gives
\begin{equation}
 \norm{\rho_{N,R}-P_R}_1\le\frac{2e}{\sqrt{1+e^2}}
 \le g(\varepsilon):=\frac{2\varepsilon}{\sqrt{1+\varepsilon^2}}\le f(\varepsilon).
 \label{eq:av1-reverse}
\end{equation}
The last inequality is $(1+\varepsilon)^2\ge1+\varepsilon^2$. Before comparing
the producers, the skeptic re-derived $g$ by a third route, purification plus
partial-trace contractivity. Neither inequality estimates a norm of an outside
product. The volume-dependent $n$ and $\varphi_{\rm out}$ cancel in every
ratio, and $\varepsilon$ depends only on anchored sums at the two sites of
$R$, which AM2 bounds independently of $N$ and of the on-site cutoff.

\subsection{The normalization trap}
\label{sec:av1-trap}
Both $\langle\psi,A\psi\rangle$ and $\norm{\psi}^2=n^2(1+e^2)$ contain every
creation in the box, and $n^2$ grows with the volume. Two shortcuts fail. The
first keeps only creations meeting $R$ in the numerator and normalizes by
$1+O(t^2)$ or by the full norm. This gives a volume-dependent value.
Normalizing by the vacuum coefficient instead gives the global vacuum overlap,
which decays with the volume. The second shortcut drops the outside creations
everywhere. The value is then volume independent but wrong, because straddling
creations contract their outside legs against $\varphi_{\rm out}$. The only
legitimate cancellation is inside the ratio $e$.

Exact three-site fixtures exhibit both failures. In the forward fixture, qubit
sites $r,o,o'$ with $R=\{r\}$ carry creations $a$ on $\{r,o\}$ (straddling),
$b$ on $\{o\}$ and $g$ on $\{o'\}$. Then
\begin{equation}
 \rho_r=\frac{1}{1+a^2+b^2}\begin{pmatrix}1+b^2&ab\\ ab&a^2\end{pmatrix},
 \qquad
 \langle\psi,P_r\psi\rangle=(1+g^2)(1+b^2).
 \label{eq:av1-fixture}
\end{equation}
The decoupled creation $g$ cancels from $\rho_r$ but not from the unnormalized
weight. At $a=\tfrac13$, $b=\tfrac12$ that weight runs through
$\tfrac54,\tfrac{29}{20},\tfrac52$ for $g=0,\tfrac25,1$, while
$\rho_r=\tfrac1{49}\bigl(\begin{smallmatrix}45&6\\6&4\end{smallmatrix}\bigr)$.
The outside creation $b$ does not cancel: the off-diagonal vanishes when $b=0$.
The straddling creation enters $e$ at first order but the off-diagonal only at
second order. In the reverse producer's own fixture, naive restriction gives
$e^2=\tfrac12$ instead of the actual $\tfrac58$. There the outside creation
\emph{increases} the excitation of $R$, so naive restriction is neither exact
nor a bound. With $k$ decoupled outside sites the local overlap stays
$\tfrac8{13}$, while the unnormalized weight is $2^k$ and the global vacuum
overlap is $\tfrac4{13}2^{1-k}$. The excitedness premise
$c_I\in\bigotimes_{x\in I}Q_x\mathcal H_x$ is also necessary: a creation with
a vacuum component gives $\langle\psi_{\rm out},\delta\rangle=-\tfrac5{16}\ne0$,
and the split \eqref{eq:av1-orth} fails.

\subsection{Derived face counts and the first-order collection}
\label{sec:av1-counts}
The first-order collection is the $k=0$ term of \eqref{eq:dict-fixed},
$c^{(1)}_M=H_M^{-1}P_MV\Omega_0$. Five facts hold for every omitted face $f$.
First, $\langle\Omega_0,W_f\Omega_0\rangle=0$. Second, $W_f\Omega_0$ carries
spin $\tfrac12$ on the four distinct links of $f$ and lies in the sector of
its owner set $M_f$, with $|M_f|\ge2$. Third, it is an eigenvector of $H_0$
with eigenvalue $24$ in normalized units, that is $3$ in units of $\alpha$.
Fourth, $\norm{W_f\Omega_0}^2=\tfrac14$. Fifth, distinct faces share at most
one link, so their vectors are orthogonal. Hence
\begin{equation}
 c^{(1)}_M=-\frac{\tau}{72}\sum_{f:\,M_f=M}W_f\Omega_0,\qquad
 \norm{c^{(1)}_M}=\frac{|\tau|\sqrt{n_M}}{144},
 \label{eq:av1-first}
\end{equation}
where $n_M$ counts the faces with owner set $M$. The amplitude is
$(\tau/3)/24$ in normalized units and $(\tau/24)/3$ in units of $\alpha$.
Mixing them would give $\tau/576$ or $\tau/9$; both are rejected controls.

The counts are derived from the I1 table (Table~\ref{tab:dict-classes}) by
translation covariance. A face of class $k$ anchored at $b$ has owner set
$b+K_k$, so it contains the factor $u$ exactly when $b=u-d$ for some
$d\in K_k$. Therefore
\begin{equation}
 \#\{f:\ u\in M_f\}=\sum_{d\in S}\#\{k:\ d\in K_k\}=21+4+8+16=49 .
 \label{eq:av1-49}
\end{equation}
Here $0$ lies in all 21 omitted classes. The offset $e_x$ lies in four: the
two $xy$ classes with $r=3$ and the two $xz$ classes with $r=3$. The offset
$e_y$ lies in eight: the four $xy$ classes with $q=1$ and the four $yz$
classes with $q=1$. The offset $e_z$ lies in all sixteen $xz$ and $yz$
classes. An owner set determines its anchor, its lexicographic minimum, and
its class shape. The 49 faces therefore fall into the 15 owner sets of
Table~\ref{tab:av1-owners}.

\begin{table}[htbp]
\centering\small
\caption{Owner sets containing a factor $u=0$, derived from the I1 table.}
\label{tab:av1-owners}
\begin{tabular}{llcc}
\toprule
Shape & Owner sets containing $0$ & $n_M$ each & Faces\\
\midrule
$\{0,e_x\}$ & $\{0,e_x\}$, $\{-e_x,0\}$ & 1 & 2\\
$\{0,e_y\}$ & $\{0,e_y\}$, $\{-e_y,0\}$ & 3 & 6\\
$\{0,e_z\}$ & $\{0,e_z\}$, $\{-e_z,0\}$ & 10 & 20\\
$\{0,e_x,e_y\}$ & three translates & 1 & 3\\
$\{0,e_x,e_z\}$ & three translates & 2 & 6\\
$\{0,e_y,e_z\}$ & three translates & 4 & 12\\
\midrule
Total & 15 & & 49\\
\bottomrule
\end{tabular}
\end{table}

The faces meeting $R$ number $49+49-16=82$. The 16 faces whose owner sets
contain both $0$ and $e_z$ are the classes containing $e_z$, anchored at $0$.
Ten faces lie inside $R$, all with owner set $\{0,e_z\}$: the six $xz$ classes
with $r\le2$ and the four $yz$ classes with $q=0$, anchored at $0$. The
original $W$ is one of them. The other 72 faces meeting $R$ have straddling
owner sets; of the 27 owner sets meeting $R$, 26 straddle. In a finite box
only whole stars are retained, so every count is at most its bulk value. The
explicit boxes $\Lambda_2$ and $\Lambda_3$ hold 64 and 216 stars (1344 and 4536
faces), with the maximum 49 reached at every bulk site. The product
$4\cdot21=84$ is used only as a labelled bound. Both producers derived the
counts before comparing them with the contract's candidates. A brute-force
fine-lattice enumeration reproduced them without the class table, and
source-mutation replays confirm that editing the table or copying a count
aborts the run.

With orthogonality inside each owner set, the first-order anchored norm and
the first-order sum over the cover obey
\begin{equation}
 \norm{c^{(1)}}_a\le\frac{|\tau|}{144}\bigl(11+2\sqrt3+3\sqrt2+2\sqrt{10}\bigr)
 \le\frac{49|\tau|}{144}=:t_1,
 \qquad
 a_1:=\sum_{M\cap R\ne\varnothing}\norm{c^{(1)}_M}\le\frac{82|\tau|}{144}.
 \label{eq:av1-t1}
\end{equation}
The admitted tier-(ii) values use the triangle forms $t_1$ and $82|\tau|/144$.
The anchored norm remains an $\ell^1$ sum over owner sets; replacing it by
$\sqrt{49}$ or $\sqrt{82}$ is the rejected root-$n$ misuse.

\subsection{Two tiers at the cap}
\label{sec:av1-tiers}
\paragraph{Tier (i): crude majorant.}
The fixed point gives $t\le J\,G(t)\le J\,G(R_a)\le28|\tau|\cdot\tfrac{148}{7}=592|\tau|$,
and only this $t$ enters $\varepsilon_{\rm i}=2t+t^2$. At the cap,
$t_{\rm i}=37/6250000$ and $\varepsilon_{\rm i}=462501369/39062500000000$.

\paragraph{Tier (ii): exact first order with a self-consistent remainder.}
Since $G$ is convex and increasing and $t\le R_a$,
\begin{equation}
 \norm{c-c^{(1)}}_a\le J\bigl(G(t)-16\bigr)\le J\,t\,G'(R_a)\le352\,J\,t,
 \qquad
 t\le t_1+352\,J\,t\ \Longrightarrow\ t\le T:=\frac{t_1}{1-352J}.
 \label{eq:av1-selfconsistent}
\end{equation}
The remainder $\varrho:=352\,J\,T=T-t_1$ is charged at its positive majorant
value and is never set to zero. Two tier mixtures are rejected: combining
$c^{(1)}$ with the crude $t$, and using $t_1$ without
\eqref{eq:av1-selfconsistent}. The forward producer uses the contract's global
form $\varepsilon_{\rm fwd}=2T+T^2$. The reverse producer bounds the
first-order sum over $R$ directly by the 82 faces and charges the remainder at
both sites, $\varepsilon_{\rm rev}=a_1+2\varrho+T^2$. Since $2T=2t_1+2\varrho$
and $a_1<2t_1$, $\varepsilon_{\rm rev}<\varepsilon_{\rm fwd}$, with ratio
$\varepsilon_{\rm fwd}/\varepsilon_{\rm rev}\approx1.19510$. The forward form
double-counts the 16 faces that contain both sites.

At $\tau=\pm10^{-8}$ the inputs are $J_0=7/25000000$, $352J_0=77/781250$,
$t_1=49/14400000000$, $T=49/14398580736$,
$\varrho=3773/11248891200000000\approx3.35\times10^{-13}$ and
$a_1=41/7200000000$. They give
\begin{equation}
 \varepsilon_{\rm fwd}=\frac{1411060914529}{207319127211110301696},\qquad
 \varepsilon_{\rm rev}=\frac{461213409663624001}{80984034066839961600000000}
 \approx5.6951153\times10^{-9}.
 \label{eq:av1-epsvalues}
\end{equation}
The resulting bounds follow. The forward values are exact rationals, the
reverse values directed rational upper bounds, and the decimals readable
summaries:
\begin{align}
 D_{\rm i}^{\rm fwd}&=\frac{36133347268157653748322}{1525878906463907516326874161}
 \approx2.36803504623\times10^{-5},\label{eq:av1-Di}\\
 D_{\rm i}^{\rm rev}&\le\frac{236800700911401877571290493559933441}{10^{40}}
 \approx2.3680070\times10^{-5},\\
 D_{\rm ii}^{\rm fwd}&=\frac{585079838465912592144137406066050}{42981220507576537932303142777593983768257}
 \approx1.36124528703\times10^{-8},\label{eq:av1-Dii}\\
 D_{\rm ii}^{\rm rev}&\le\frac{113902305553947264976096174312887}{10^{40}}
 \approx1.1390231\times10^{-8}.\label{eq:av1-Diirev}
\end{align}
Every value is the same at both signs, because only $|\tau|$ enters.
Table~\ref{tab:av1-tiers} compares them with the targets.

\begin{table}[htbp]
\centering\small
\caption{AV1 state bounds at $\tau=\pm10^{-8}$ against the preregistered target
$4/10^7$ and the secondary comparison $10^{-6}$.}
\label{tab:av1-tiers}
\begin{tabular}{lrcc}
\toprule
Bound on $\norm{\rho_R-P_R}_1$ & Approximate value & $\le4/10^7$ & $\le10^{-6}$\\
\midrule
Round31 reset, $2\sqrt{49|\tau|/3}$ & $8.08290376865\times10^{-4}$ & no & no\\
Tier (i), forward exact & $2.36803504623\times10^{-5}$ & no & no\\
Tier (i), reverse directed & $2.3680070\times10^{-5}$ & no & no\\
Tier (ii), forward exact (both inequalities) & $1.36124528703\times10^{-8}$ & yes & yes\\
Tier (ii), reverse directed (fidelity) & $1.1390231\times10^{-8}$ & yes & yes\\
\bottomrule
\end{tabular}
\end{table}

Tier (i) fails both targets. It is retained as a limited-tier value and is not
retuned. Iterating the fixed-point bound once (forward) or four times
(reverse) lowers it only to about $1.79\times10^{-5}$. Tier (ii) meets the
target with margins $29.38$ (forward) and $35.12$ (reverse). The two tier-(ii)
numbers have different backing. $D_{\rm ii}^{\rm fwd}$ is certified by both
inequalities, because $g(\varepsilon_{\rm fwd})\le f(\varepsilon_{\rm fwd})$.
$D_{\rm ii}^{\rm rev}$ needs the fidelity inequality: the forward formula
evaluated at $\varepsilon_{\rm rev}$ exceeds it. Both are valid upper bounds
on the same quantity. The gate binds $D_{\rm ii}^{\rm fwd}$ as the state
premise for AV2.

\subsection{From the cutoff box to the AQ state}
\label{sec:av1-passage}
\paragraph{Cutoff-vector removal.}
AM2 removed the on-site cutoff for the eigenvalues and the gap only; a reduced
density needs the ground vector. Fix $N$ and let
$Q_L=\bigotimes_{x\in\Lambda_N}1_{[0,L]}(h_x)$. This is a finite-rank
projection that commutes with $H_0$ and increases strongly to $I$. Let
$\psi_L$ be the normalized ground vector of $Q_LHQ_L$. AM2 applies in each
cutoff space with rank-independent constants. For $L\ge24$,
$Q_LW_f\Omega_0=W_f\Omega_0$ and $c^{(1)}_L=c^{(1)}$; for smaller $L$ each
face vector is kept or annihilated, so the counts only decrease. Hence
$\norm{\rho_{L,R}-P_R}_1\le D$ for every $L$. Monotone convergence of
$\langle Q_L\psi,H_0Q_L\psi\rangle$ and boundedness of $V$ in the fixed box
give $E_0\le E_{0,L}\to E_0$. The untruncated gap $E_1-E_0\ge\tfrac12$ (AM2
\S6) then gives the Eckart-type bound
\begin{equation}
 1-|\langle\psi,\psi_L\rangle|^2\le\frac{E_{0,L}-E_0}{E_1-E_0}
 \le2\,(E_{0,L}-E_0)\longrightarrow0 .
 \label{eq:av1-eckart}
\end{equation}
The pure-state trace distance is $2\sqrt{1-|\langle\psi,\psi_L\rangle|^2}$,
and partial trace is trace-norm contractive, so $\rho_{L,R}\to\rho_{N,R}$. The
closed trace-norm ball keeps $\norm{\rho_{N,R}-P_R}_1\le D$ for the untruncated
finite-volume ground state. The reverse route reaches the same conclusion from
a Rayleigh quotient with the uniform cutoff gap $\tfrac12$. The gap is
essential, because eigenvalue convergence alone does not control vectors. In
$\mathrm{diag}(0,0,1)$ the cutoff ``ground'' $e_2$ has the exact ground energy
yet is orthogonal to $e_1$; the alternating family $\mathrm{diag}(0,1/n)$,
$\mathrm{diag}(1/n,0)$ has converging eigenvalues and alternating ground
vectors.

\paragraph{Passage to the AQ state.}
AQ1 extracts subsequences $N_k$ along which every local density converges in
trace norm. For any such subsequence with limit density $\rho_R$,
\begin{equation}
 \norm{\rho_R-P_R}_1\le\norm{\rho_R-\rho_{N_k,R}}_1+D\longrightarrow D .
 \label{eq:av1-passage}
\end{equation}
The quantifier is \emph{every} subsequential limit of AQ1's centered
whole-star construction at either sign, AQ1's chosen state included. Two such
limits lie within $2D$ of each other on $R$, and the checker exhibits two
distinct limits inside the same ball. This is uniform local closeness, not
uniqueness, whole-sequence convergence or a rate in $N$. Trace-norm
convergence is needed: a weak-* limit can lose mass, as
$(1-D/2)|e_0\rangle\langle e_0|+(D/2)|e_n\rangle\langle e_n|$ shows, and
AQ1's energy tightness excludes that loss.

\subsection{Consequences, scaling and AV2 feasibility}
\label{sec:av1-consequences}
Trace duality with $\norm{W}\le1$ bounds the mean. For the second moment, a
trace-zero difference has expectation at most half its trace norm against an
effect, and $0\le W^2\le I$. Hence
\begin{equation}
 |\omega_\tau(W)|\le D,\qquad
 \bigl|\omega_\tau(W^2)-\tfrac14\bigr|\le\tfrac D2,\qquad
 m_\tau^2\le D^2 .
 \label{eq:av1-consequences}
\end{equation}
At tier (ii) these bounds are about $1.3612\times10^{-8}$ and
$6.806\times10^{-9}$. The mean-square charge $m_\tau^2\le D^2$ is kept in
every later budget. Neither the value nor the sign of $\omega_\tau(W)$ is
claimed. A candidate first-order mean $|\tau|/144\approx6.94\times10^{-11}$ is
consistent with the bound but is neither implied nor claimed; it belongs to
AW1.

Every bound is increasing in $|\tau|$, so the cap values bound every smaller
coupling. The preregistered scaling test compares the cap with $10^{-10}$:
\begin{equation}
 \frac{D(10^{-8})}{D(10^{-10})}\approx
 \begin{cases}
 100.0098 & \text{tier (ii), forward},\\
 100.0117 & \text{tier (ii), reverse},\\
 100.0015 & \text{tier (i), forward},\\
 100.0003 & \text{tier (i), reverse},
 \end{cases}
 \qquad
 \frac{2\sqrt{49\cdot10^{-8}/3}}{2\sqrt{49\cdot10^{-10}/3}}=10 .
 \label{eq:av1-scaling}
\end{equation}
All tier ratios lie in $[99,101]$, so the bounds are linear in $\tau$. The
Round31 bound has ratio exactly 10, and relabelling it as linear is a rejected
control. At the cap, tier (ii) is about $5.9\times10^{4}$ times smaller than
the Round31 reset bound.

With a directed Machin bracket for $\pi$, the AV2 feasibility function
$F(D)=2(D+D^2)+49\cdot10^{-8}/\pi$ satisfies
$F(D_{\rm ii}^{\rm fwd})\le1.8320\times10^{-7}$,
$F(D_{\rm ii}^{\rm rev})\le1.7876\times10^{-7}$ and
$F(4/10^7)\le9.5597\times10^{-7}$. It meets $10^{-6}$ at $4.22\times10^{-7}$
and fails at $4.23\times10^{-7}$ and at $D_{\rm i}$. This is arithmetic on the
contract statement only; AV1 does not prove the window lemma.

\subsection{Error ledger, controls and replays}
\label{sec:av1-evidence}
The forward producer exports the preregistered error terms at tier (ii).
\begin{itemize}
\item \recid{am2_remainder}: $\varrho\approx3.35\times10^{-13}$; not
applicable to tier (i), which bounds the whole collection directly.
\item \recid{two_creation}: $T^2\approx1.16\times10^{-17}$, inside
$\varepsilon$.
\item \recid{straddling}: inside the $2t$ term; at first order the 72
straddling faces contribute at most $|\tau|/2$ to the sum over $R$.
\item \recid{density}: the $\sigma$ and $d^2P_R$ part of $D$,
$2\varepsilon^2/(1+\varepsilon^2)\approx9.3\times10^{-17}$.
\item \recid{onsite_cutoff_vector}: no numeric cost; it is an exact
per-box limit.
\item \recid{arithmetic}: no numeric cost; every bound is an exact
rational, and the only enclosures are directed upward
($e^{1/8}<8/7$, $e^{8t}\le1/(1-8t)$, integer square-root brackets).
\end{itemize}
The reverse producer exports no six-name ledger. The skeptic's mapping
supplies it: \recid{am2_remainder} is $2\varrho$, \recid{density} is not
applicable to the fidelity route, and \recid{arithmetic} is outward rounding
of at most $10^{-40}$ on $D$.

The contract listed 25 controls. The reverse checker, with 54 exact checks,
implements all 25 as damaging mutations that must raise explicit exceptions.
The forward checker, with 43 checks, implements 22 of them as mutations. The
other three are positive exact checks that match their prospective semantics:
\recid{no_priority_or_continuum_claim},
\recid{first_order_face_enumeration} and
\recid{av2_feasibility_threshold}. The forward report's sentence claiming
that every control is a mutation is therefore not admitted. The skeptic's
pre-comparison package (71 exact checks) was hash-recorded before either
producer was opened. The post-comparison review adds 39 checks, including 12
source-mutation replays. Both producers reproduce their frozen outputs byte
for byte under normal and optimized Python. The reverse input inventory is
exactly \texttt{AGENTS.md}, the contract and its shared premises, as enforced
by the freezing tool.

\subsection{Limitations recorded by the gate}
\label{sec:av1-limits}
The AV1 gate records the following limitations verbatim.
\begin{gatequote}
\begin{itemize}
\item Only the AM2/AQ1 zero-selected patterned family, the cover R=\{0,e\_z\}
and fixed spacing; nothing transfers to nonzero selected triples (non-Haar
P\_R), uniform Wilson theory, weak coupling or the continuum
\item Upper certificates only: no lower bound on ||rho\_R-P\_R||\_1; value and
sign of omega(W), the first-order coefficient 1/144 and any first-order
vanishing are not claimed (AW1)
\item Uniform local closeness of all AQ1 subsequential limits (pairwise within
2D on R), not uniqueness, whole-sequence convergence or a rate in N
\item Inherited without re-proof: AM2 fixed point, multilinear majorant,
cutoff-uniform and untruncated gap 1/2 (AM2 section 6); AQ1 trace-norm
compactness and diagonal extraction; the I1 dictionary
\item Product-ordering split and normalization-trap location are shared panel
premises; independence limited to the inequality, constants, enumeration,
cutoff and passage steps; all agents are correlated model agents
\item Sign-blind: the tau=-10\textasciicircum{}-8 evaluation replays the same
|tau| formula
\item Tier (i) fails 4/10\textasciicircum{}7 and 10\textasciicircum{}-6 and is
retained as a limited-tier value; AT4 square-root bound not used
\item D\_ii(rev)\textasciitilde{}1.1390e-8 rests on the fidelity inequality
and the labelled 82-face refinement; D\_ii(fwd)\textasciitilde{}1.3612e-8 is
certified by both inequalities
\item Forward check.py implements 22 of 25 contract controls as
exception-raising mutations; no\_priority\_or\_continuum\_claim,
first\_order\_face\_enumeration and av2\_feasibility\_threshold are positive
exact checks; the forward sentence "Every control is a damaging mutation that
must raise an explicit exception." is not admitted
\item Reverse exports no six-name error ledger; the skeptic mapping in av1.md
(N2) supplies it
\item Reverse remark "no parity argument removes a first-order mean" is not
admitted (AW1 exclusion)
\item Forward attribution of a "2(E\_\{0,L\}-E\_0)/gap form" to the contract
is incorrect wording; the Eckart inequality itself is correct
\item Both producers disclose undeclared protocol/code-style reads (tools
README and freeze.py, AM2 and AT5 check.py); none carries premise weight
\item Scientific priority unverified
\end{itemize}
\end{gatequote}
Full derivations and reviewed implementations:
\repo{research/round32/forward/av1/report.md}{AV1 forward},
\repo{research/round32/reverse/av1/report.md}{AV1 reverse},
\repo{research/round32/skeptic/av1.md}{independent model-agent review} and the
\repo{research/round32/contracts/av1.json}{frozen contract}. AV1 by itself
certifies no Euclidean datum; Section~\ref{sec:av2} uses its forward tier-(ii)
value.
